\chapter{Text I/O: a pure-Ada Ada.Text\_IO over the console and ext4}
\label{ch:text_io}

The standard library's \texttt{Ada.Text\_IO} is how an Ada programmer reaches the
outside world: \texttt{Put\_Line}, formatted numbers, files. On this bare runtime
that package is \emph{stripped}: it prints characters to the console and reads one
back, nothing more. There are no files, no \texttt{Integer\_IO}, no layout.
Chapter~\ref{ch:runtime-build} explains why. The full package will not fit a
configurable run-time.

\texttt{ESP32S3.Text\_IO} fills the gap. It is a near-complete, \emph{pure-Ada}
re-implementation of \texttt{Ada.Text\_IO}. It works over two things at once: the
console (the USB-Serial-JTAG port, through the device mux of
Section~\ref{driver:console-mux}) and \textbf{ext4 files} on a mounted volume
(Chapter~\ref{ch:storage}). There is no C anywhere in it. In particular there is
no \texttt{Interfaces.C\_Streams}, the C \texttt{FILE*} layer the hosted GNAT
\texttt{Ada.Text\_IO} is built on. Files are reached through the ext4 VFS, and
numbers are formatted by hand in Ada.

This chapter shows how to use it, then explains the one thing it cannot do:
\emph{be} \texttt{Ada.Text\_IO}.

\section{Console output}\index{ESP32S3.Text\_IO}

The file-less subprograms write to the current output file, which starts as the
console:

\begin{lstlisting}
with ESP32S3.Text_IO; use ESP32S3.Text_IO;
...
Put_Line ("Lua on the ESP32-S3");
Put ("count = ");  Put (Integer'Image (N));  New_Line;
\end{lstlisting}

\cnote{this is the role \texttt{printf} and \texttt{ESP\_LOGI} play. For terse
status logging the lighter \texttt{ESP32S3.Log} is often enough
(Chapter~\ref{ch:hal}). Reach for \texttt{ESP32S3.Text\_IO} when you want files,
the numeric generics, or column layout.}

The console is a \texttt{File\_Type} like any other. \texttt{Standard\_Output},
\texttt{Standard\_Error} and \texttt{Standard\_Input} name it, and
\texttt{Set\_Output} redirects the file-less calls:

\begin{lstlisting}
Put_Line (Standard_Output.all, "explicit form, same console");
Set_Output (Some_File_Access);    -- now file-less Put_Line goes to that file
\end{lstlisting}

\subsection{Three destinations: USB, a UART, a file}\index{Set\_Output}

Two knobs choose where text goes. \texttt{ESP32S3.Serial.Set\_Output} picks the
physical console device, and everything console-bound follows it:
\texttt{ESP32S3.Log}, the runtime's own \texttt{Ada.Text\_IO}, and this package's
console writes. \texttt{ESP32S3.Text\_IO.Set\_Output} picks where this package's
file-less calls land: the console, or a file.

\begin{lstlisting}
--  1. The built-in USB Serial/JTAG console (the default).
ESP32S3.Serial.Set_Output (ESP32S3.Serial.Console_Device);

--  2. A UART at 115_200 baud, 8-N-1, transmit on IO17.
Port : aliased ESP32S3.UART.Session;        -- must outlive the redirect
...
ESP32S3.UART.Acquire (Port, ESP32S3.UART.UART1, Baud => 115_200, Tx => 17);
ESP32S3.Serial.Set_Output (ESP32S3.UART.Text.As_Device (Port));

--  3. A file on the ext4 volume.
Log : ESP32S3.Text_IO.File_Type;
...
ESP32S3.Text_IO.Create     (Log, "/flash/console.log");
ESP32S3.Text_IO.Set_Output (Log);
\end{lstlisting}

Cases 1 and 2 move the whole console at the device level
(Section~\ref{driver:console-mux}). Case 3 sends only the file-less output of
\texttt{ESP32S3.Text\_IO} to the file because a UART and the USB port have no
file behind them: a file is reached through a \texttt{File\_Type}, not the device
mux. Point output back at the console with \texttt{Set\_Output (Standard\_Output.all)}.

\section{Files on ext4}\index{ext4 files}\index{File\_Type}

A file is named by a VFS path. The first component selects the mount
(\texttt{/flash}, \texttt{/sd}), and the rest is the path within it. Open or
create, read or write lines, then close:

\begin{lstlisting}
declare
   F : File_Type;
begin
   Create (F, "/flash/log.txt");           -- Out_File by default
   Put_Line (F, "boot ok");
   Put_Line (F, "temp = 23.4 C");
   Close (F);                               -- commits to flash

   Open (F, "/flash/log.txt", In_File);
   while not End_Of_File (F) loop
      Put_Line (Get_Line (F));              -- echo the file to the console
   end loop;
   Close (F);
end;
\end{lstlisting}

\texttt{File\_Type} is controlled. A file left open is closed (and committed) when
it goes out of scope, so a forgotten \texttt{Close} loses nothing. Append mode and
the rest of the management surface work as in the Reference Manual: \texttt{Reset}
(reposition, or re-open with a new mode), \texttt{Mode}, \texttt{Name},
\texttt{Flush}, and \texttt{Delete}.

\textbf{No store, no crash.} Errors are ordinary Ada exceptions, because the ext4
layer's exceptions \emph{are} renames of \texttt{Ada.IO\_Exceptions}. With no
volume mounted, opening a path on it raises \texttt{Name\_Error}. A read past the
end raises \texttt{End\_Error}, and the wrong mode raises \texttt{Mode\_Error}. So
a program that logs to flash degrades to a catchable exception, never a fault:

\begin{lstlisting}
begin
   Open (F, "/flash/log.txt", In_File);
   ...
exception
   when Name_Error => null;    -- no flash mounted: skip the file, keep running
end;
\end{lstlisting}

\textbf{Durability.} By default a file commits on \texttt{Close}. The
implementation-defined \texttt{Form} string carries one ext4 option,
\texttt{"sync=yes"}, which commits after \emph{every} write. That is crash-durable,
at the cost of a journal flush per line. An unrecognised form raises
\texttt{Use\_Error}.

\begin{lstlisting}
Create (F, "/flash/journal.txt", Form => "sync=yes");
\end{lstlisting}

\section{Formatted output: the generics}\index{Integer\_IO}\index{Float\_IO}

The full family of numeric and enumeration generics is present. Each has
\texttt{Put} \emph{and} \texttt{Get}, the \texttt{To} and \texttt{From}
\texttt{String} variants, and defaults from the type's attributes:

\begin{lstlisting}
package Int_IO is new Integer_IO (Integer);
package Flt_IO is new Float_IO   (Float);

Int_IO.Put (255, Width => 8, Base => 16);   --  "  16#FF#"
Flt_IO.Put (3.14159, Fore => 1, Aft => 2);  --  "3.14"

declare
   N : Integer;  Last : Positive;
begin
   Int_IO.Get ("  -42 rest", N, Last);       --  N = -42, Last = 5
end;
\end{lstlisting}

\texttt{Integer\_IO}, \texttt{Modular\_IO}, \texttt{Float\_IO}, \texttt{Fixed\_IO},
\texttt{Decimal\_IO} and \texttt{Enumeration\_IO} are all there. Floating output is
fixed-point, or scientific when \texttt{Exp > 0}. The fixed-point generics format
through a wide intermediate, so a ranged \texttt{delta} type cannot overflow while
it is being scaled. None of it uses a host \texttt{sprintf}. The digits are
produced in Ada.

\section{Column and line layout}\index{Set\_Col}\index{column layout}

The same line, column and page model the Reference Manual defines is here, which
is how you align columns without counting spaces:

\begin{lstlisting}
Put (F, "voltage"); Set_Col (F, 16); Put_Line (F, "3.30");
Put (F, "id");      Set_Col (F, 16); Put_Line (F, "42");
--  voltage         3.30
--  id              42
\end{lstlisting}

\texttt{Col}, \texttt{Line}, \texttt{Page}, \texttt{Set\_Col}, \texttt{Set\_Line},
\texttt{New\_Page} and \texttt{Skip\_Page} behave as in the RM. Setting a non-zero
\texttt{Line\_Length} turns on automatic wrapping. Like the RM, the position only
ever moves forward. A stream cannot un-write a byte, so \texttt{Set\_Col} to an
earlier column starts a new line rather than seeking back.

\section{Why this is not Ada.Text\_IO}\index{configurable run-time}

The obvious question is why all of this is named \texttt{ESP32S3.Text\_IO} instead
of simply \emph{being} \texttt{Ada.Text\_IO}, so that legacy \texttt{with
Ada.Text\_IO} code would use it unchanged. We tried. The compiler will not allow
it.

\texttt{Ada.Text\_IO} is special to GNAT. Its nested numeric generics
(\texttt{Integer\_IO}, \texttt{Float\_IO}, and the rest) are not ordinary generics
that the compiler reads from your source. GNAT recognises the unit by name and
binds those generics to its \emph{own} predefined implementation units
(\texttt{a-tiinio}, \texttt{a-tiflio}, and so on). You cannot supply your own.

Worse, those predefined units are exactly what a \textbf{configurable run-time}
forbids. Build a full \texttt{Ada.Text\_IO} spec into this runtime and the compiler
stops at the first generic:

\begin{lstlisting}[style=sh]
a-textio.ads: error: construct not allowed in configurable run-time mode
a-textio.ads: error: file a-tiinio.ads not found  ("Integer_Io" not available)
\end{lstlisting}

That is the whole reason the bare runtime ships a \emph{stripped}
\texttt{Ada.Text\_IO} in the first place. The standard one, with its formatted
I/O, depends on machinery that a configurable run-time excludes. The console hook
of Section~\ref{driver:console-mux} does route the stripped package's
\texttt{Put\_Line} through our device mux, so console-only legacy code already
benefits. But the files and the generics cannot live under the
\texttt{Ada.Text\_IO} name on this target.

So \texttt{ESP32S3.Text\_IO} stands beside \texttt{Ada.Text\_IO}, not in place of
it. The cost is a one-line change to legacy source: \texttt{with} and \texttt{use}
the new name. The benefit is a Reference-Manual-shaped, pure-Ada text I/O over the
console and ext4, with formatted numbers and column layout, on a runtime the
standard package was never going to fit.
